\begin{exercise}
\label{ex:first-incompleteness-01}
In the proof sketch of \cref{thm:godel-first}, identify the precise uses of
effectiveness, representability, consistency, and \(\Sigma^0_1\)-soundness.
\end{exercise}

\begin{exercise}
\label{ex:first-incompleteness-02}
Show that if \(T\) is consistent, then the G\"odel sentence \(G_T\) constructed
above is true in \(\mathcal N\).  Explain why this does not imply that
\(T\cup\{G_T\}\) proves every arithmetical truth.
\end{exercise}

\begin{exercise}
\label{ex:first-incompleteness-03}
Why can the proof sketch of Rosser's theorem compare proof codes with a
bounded
quantifier without assuming omega-consistency?  State the role of the least
proof code.
\end{exercise}

